\section{Action-Grounded World Pretraining}
\label{sec:pretrain}

\sref{sec:joint} states the objective and the measurement instruments it
implies. This section gives the realisation: the two streams and the geometry
that lets them share one attention space, the interfaces that let heterogeneous
sources enter one objective, the corpus, the pretraining procedure, and the
mapping from the formulation to the measurements of \sref{sec:exp}.

\subsection{Two Streams in One Attention Space}
\label{sec:streams}

\Paragraph{World stream.}
The world stream is a pretrained text-and-image-to-video diffusion transformer of
five billion parameters from the Wan family \citep{wan2025wan}, used with its
official weights and its official implementation. It carries thirty layers of
width 3072 arranged as twenty-four attention heads of width 128. Its frozen variational
autoencoder \citep{kingma2014vae,rombach2022ldm} compresses space by sixteen and
time by four into forty-eight latent channels, and its frozen text encoder supplies a variable-length language
context. Neither frozen module holds gradient, and both are evaluated offline
(\sref{sec:procedure}).

\Paragraph{Action stream.}
The action stream is a pre-norm transformer of thirty layers and width 1024 with
adaptive layer normalisation driven by its own noise level, one-dimensional
rotary positions over the chunk, and its own text cross-attention projections.
Its readout heads are zero-initialised, so training starts from a model whose
world behaviour equals the video prior and whose action output is null. The
stream projects queries, keys, and values into twenty-four heads of width 128,
matching the world backbone exactly, which is the condition that lets the two
streams share one attention space while keeping different residual widths. The
construction is a mixture of two experts over one attention operator, applied
across the action and world modalities.

\Paragraph{The exchange.}
World layer $i$ pairs with action layer $j$. The action block projects its
queries, keys, and values first and publishes the keys and values. The world
block then runs with an attention processor that appends those keys and values
when the visibility setting lets the world read the action. The action block
finally attends over the concatenation of the world keys and values with its
own, optionally through a stop-gradient, and finishes its block. Attention
carries no parameters, so the coupling costs exactly the parameters of the
second stream and nothing else. One boolean mask built from the coupling
specification decides every visibility question, and no other part of the
forward pass reads the coupling.

\Paragraph{Token layout.}
The action-side sequence is the state condition token, then thirty-two action
tokens, then sixteen consequence tokens when the consequence stream is enabled,
which gives forty-nine action-side tokens. A window samples nine observation
frames at stride four across thirty-two raw control steps, and the autoencoder
turns them into three latent frames of which the first is the clean present
observation. The observation canvas of \sref{sec:interfaces} yields 288 world
tokens for a two-camera platform and 360 for a three-camera platform. The action
side is therefore a small fraction of the joint sequence, and the cost of the
coupling is dominated by the second stream rather than by the attention.

\subsection{Interfaces That Admit Heterogeneous Sources}
\label{sec:interfaces}

\Paragraph{Observation canvas.}
Every platform composes its cameras into one canvas whose slots are declared in
configuration, so multiple views enter as a single video stream and the token
count stays independent of the camera count. Each view keeps its aspect ratio
inside its slot, absent views are zero-filled, and canvas dimensions are
multiples of thirty-two. A two-camera platform uses a canvas of height 384 and
width 256 with a 256 by 256 head view above a 128 by 128 wrist view. A
three-camera bimanual platform uses height 384 and width 320 with a 256 by 320
head view above two 128 by 160 wrist views. \fref{fig:interfaces} draws both the
canvas and the two vector layouts.

\Paragraph{Action layout.}
All sources write into one typed layout of width eighty. Each arm contributes an
end-effector position of three dimensions, a rotation in the six-dimensional
continuous representation, one gripper aperture, and twenty-four hand joints,
which gives thirty-four dimensions per arm and sixty-eight for two arms; twelve
slots stay reserved. A source declares a binding that routes each range of its
raw vector to a named slot, and the layout derives the per-dimension validity
mask from that binding. The mask enters the action term of \eqref{eq:loss}
directly, so a single-arm platform contributes gradient to the left-arm slots
alone and a source without action labels contributes to the world term alone. No
index arithmetic lives outside the layout object, which is what makes the
binding checkable rather than conventional.

\Paragraph{Normalisation.}
Statistics are computed per source over its training split and frozen into the
cache fingerprint. Positions and joints receive quantile normalisation into
$[-1,1]$. The rotation slots are pinned to the identity range so the continuous
rotation representation passes through unchanged. The gripper slot receives a
fixed affine map that sends a closed gripper to the low end and an open gripper
to the high end, which removes a convention that differs between sources.
Deployment inverts the same map before commanding the robot.

\Paragraph{Consequence layout.}
A consequence label uses forty-eight global keypoint slots over a horizon of
sixteen control steps with the displacement gain of \eqref{eq:consequence} set
to four. Seven slots track keypoints of the manipulator in the head view and
seven track them in the wrist view, each anchored at its own position at the
observation step. Twenty-five slots form a five by five probe grid over the
wrist view anchored at fixed query locations, which gives coverage of the scene
rather than the actuator alone. Labels come from an off-the-shelf point tracker
run offline \citep{doersch2023tapir,karaev2024cotracker} and are stored as
sidecar files beside the episodes, so the consequence stream adds no cost to the
training loop.

\begin{figure}[t]
\centering
\adjustbox{max width=\linewidth}{%
\begin{tikzpicture}[font=\figfont, x=1mm, y=1mm,
  bx/.style={draw=black!65, line width=0.5pt},
  lbl/.style={align=center},
  ttl/.style={anchor=west}]

% ---- (a) observation canvas ---- %
\node[ttl] at (0,52) {\textbf{(a)} observation canvas, one video stream per platform};
\filldraw[fill=moSage!45, bx] (0,14) rectangle (48,42.4);
\filldraw[fill=moSage!25, bx] (0,0) rectangle (24,14);
\filldraw[fill=moStone!30, bx] (24,0) rectangle (48,14);
\node[lbl] at (24,28.2) {head view, $256\times320$};
\node[lbl] at (12,7) {left wrist\\$128\times160$};
\node[lbl] at (36,7) {absent view\\zero-filled};
\draw[{Stealth[length=2mm]}-{Stealth[length=2mm]}, black!65] (-5,0) -- (-5,42.4)
  node[midway, anchor=south, rotate=90] {384};
\draw[{Stealth[length=2mm]}-{Stealth[length=2mm]}, black!65] (0,-5) -- (48,-5)
  node[midway, anchor=north] {320};
\node[lbl, anchor=west, align=left] at (56,21)
  {aspect ratio preserved inside each slot;\\
   canvas sides are multiples of thirty-two;\\
   the token count stays independent of\\ the camera count};

% ---- (b) action layout ---- %
\begin{scope}[shift={(0,-34)}]
  \node[ttl] at (0,20) {\textbf{(b)} shared action layout, width 80};
  \filldraw[fill=moBlue!50, bx] (0,0) rectangle (68,11);
  \filldraw[fill=moClay!55, bx] (68,0) rectangle (136,11);
  \filldraw[fill=moStone!30, bx] (136,0) rectangle (160,11);
  \foreach \x in {6,18,20} {\draw[black!50, line width=0.4pt, densely dotted] (\x,0) -- (\x,11);}
  \foreach \x in {74,86,88} {\draw[black!50, line width=0.4pt, densely dotted] (\x,0) -- (\x,11);}
  \draw[bx] (0,0) rectangle (160,11);
  \node[lbl] at (34,5.5) {left arm, 34 slots};
  \node[lbl] at (102,5.5) {right arm, 34 slots};
  \node[lbl] at (148,5.5) {reserved 12};
  \foreach \x/\t in {0/0, 68/34, 136/68, 160/80} {
    \draw[black!65, line width=0.5pt] (\x,11) -- (\x,14.5);
    \node[anchor=south] at (\x,14.7) {\t};}
  \node[lbl, anchor=north west, align=left] at (0,-3)
    {each arm in order: end-effector position 3, rotation 6, gripper 1, hand joints 24;
     the dotted lines mark those internal boundaries and the routing generates the validity mask};
\end{scope}

% ---- (c) consequence layout ---- %
\begin{scope}[shift={(0,-72)}]
  \node[ttl] at (0,20) {\textbf{(c)} consequence slots, 48 over a horizon of 16 control steps};
  \filldraw[fill=moRose!60, bx] (0,0) rectangle (23.3,11);
  \filldraw[fill=moRose!38, bx] (23.3,0) rectangle (46.6,11);
  \filldraw[fill=moRose!18, bx] (46.6,0) rectangle (130,11);
  \filldraw[fill=moStone!30, bx] (130,0) rectangle (160,11);
  \draw[bx] (0,0) rectangle (160,11);
  \node[lbl] at (11.6,5.5) {head\\keypoints 7};
  \node[lbl] at (35,5.5) {wrist\\keypoints 7};
  \node[lbl] at (88,5.5) {wrist probe grid, 25 slots on a $5\times5$ pattern};
  \node[lbl] at (145,5.5) {unused 9};
  \foreach \x/\t in {0/0, 23.3/7, 46.6/14, 130/39, 160/48} {
    \draw[black!65, line width=0.5pt] (\x,11) -- (\x,14.5);
    \node[anchor=south] at (\x,14.7) {\t};}
  \node[lbl, anchor=north west, align=left] at (0,-3)
    {each slot carries a displacement and an observability value per step;
     the fourteen actuator slots are anchored at their tracked position and the grid slots at fixed queries};
\end{scope}
\end{tikzpicture}}
\caption{The three interfaces that let heterogeneous sources enter one
objective. (a) Cameras compose into one canvas at a fixed token count. (b) Every
source routes its raw action vector into named slots of one layout of width
eighty, and the routing generates the per-dimension validity mask that the action
term uses, so a single-arm source writes into the left block alone. (c) A
consequence label carries the displacement and the observability of each
keypoint slot over sixteen control steps.}
\label{fig:interfaces}
\end{figure}

\subsection{The Pretraining Corpus}
\label{sec:corpus}

The corpus mixes two kinds of source. Robot sources carry action labels and
supervise both terms of \eqref{eq:loss}. Egocentric human sources carry no
action label, so their action mask is empty and they supervise the world term
and, where a tracker has produced labels, the consequence term.
\tref{tab:corpus} lists every source with the quantities the recipe file records
and the code checks; the robot sources are read through a shared episode format
\citep{cadene2024lerobot} where one exists. The recipe reaches 6\,451 hours, which slightly exceeds the
6\,369 hours of the largest published world-action recipe
\citep{referencestack}. Hours are computed from frame counts divided by verified
control rates rather than trusted from container metadata, because several
public releases carry incorrect rate labels; the two to six times discrepancies
we found are recorded per source in the recipe file. Mixture weights follow a
temperature rule $p_i \propto n_i^{0.7}$ over window counts, which keeps a
source with many short episodes from dominating a source with fewer long ones.

The table separates two tiers, and the separation is an availability statement
rather than a scientific one. The first tier is already on local storage and
reads through existing adapters, which lets a run start without any transfer.
The second tier is staged from public releases; the largest single item is a
bimanual simulation corpus that the reference recipe also used, and the
remainder adds egocentric diversity and two further real platforms. Both tiers
enter the same objective through the interfaces of \sref{sec:interfaces}, so a
partial corpus simply reduces the hour budget of an arm.

Two exclusions define the evaluation discipline. The benchmark corpora of
\sref{sec:setup} stay out of pretraining entirely, so an out-of-distribution
claim never rests on a source that shares its distribution with the evaluation
suite; this is what removes the large simulation share that the reference recipe
carried. One real source with a three hertz control rate is retained with a
shortened window, because a thirty-two step chunk spans more than ten seconds at
that rate and would otherwise mix time scales with the two to three second
windows of every other source.

\begin{table}[t]
\centering
\small
\caption{The pretraining corpus. Hours are frame counts divided by verified
control rates. Robot sources supervise the world term and the action term;
egocentric human sources supervise the world term and, where tracker labels
exist, the consequence term. Tier~1 is on local storage and reads through
existing adapters; tier~2 is staged from public releases. The total slightly
exceeds the 6\,369 hours of the largest published world-action recipe
\citep{referencestack}, so pretraining comparisons in this report are made at a
comparable scale and differ from that recipe in composition rather than in
size. The total divides into 3\,933 hours of real robot data, 146 hours of
simulation, and 2\,372 hours of egocentric human footage.}
\label{tab:corpus}
\adjustbox{max width=\linewidth}{%
\begin{tabular}{@{}llrrcc@{}}
\toprule
\textbf{Source} & \textbf{Platform} & \textbf{Episodes} & \textbf{Hours} & \textbf{Action} & \textbf{Consequence} \\
\midrule
\panel{6}{Tier 1, real robots with action labels} \\
DROID \citep{khazatsky2024droid} & Franka, three views & \phantom{0}95\,658 & \phantom{0}440 & yes & partial \\
RT-1 \citep{brohan2023rt1} & Google Robot & \phantom{0}86\,403 & \phantom{0}350 & yes & no \\
RoboMIND, on disk & Franka, AgileX, UR & \phantom{0}\phantom{0}\phantom{0}n/a & \phantom{0}130 & yes & no \\
BridgeData V2 \citep{walke2023bridgedata} & WidowX-250 & \phantom{0}53\,192 & \phantom{0}111 & yes & no \\
Language-Table & planar arm & 179\,000 & \phantom{00}80 & yes & no \\
RoboSet & Franka, multi-camera & \phantom{0}\phantom{0}\phantom{0}n/a & \phantom{00}22 & yes & no \\
RH20T & UR5, KUKA, Flexiv & \phantom{00}5\,474 & \phantom{00}20 & yes & no \\
\addlinespace[2pt]
\panel{6}{Tier 1, simulation outside the evaluation suites} \\
RoboCasa GR1 & humanoid, dexterous hands & \phantom{0}24\,000 & \phantom{0}118 & yes & no \\
CALVIN \citep{mees2022calvin} & Franka & \phantom{0}18\,957 & \phantom{00}15 & yes & partial \\
Other small suites & mixed & \phantom{0}\phantom{0}\phantom{0}n/a & \phantom{00}13 & yes & no \\
\addlinespace[2pt]
\panel{6}{Tier 1, egocentric human without action labels} \\
EgoDex & two human hands & 318\,082 & \phantom{0}775 & no & yes \\
Ego4D \citep{grauman2022ego4d} & unconstrained & \phantom{00}\phantom{0}665 & \phantom{0}295 & no & no \\
\addlinespace[3pt]
\panel{6}{Tier 2, staged from public releases} \\
Bimanual simulation corpus & bimanual, two views & \phantom{0}\phantom{0}\phantom{0}n/a & 1\,900 & yes & no \\
Egocentric remainder & unconstrained & \phantom{0}\phantom{0}\phantom{0}n/a & 1\,180 & no & no \\
Large-scale bimanual real & dual-arm, two wrists & \phantom{0}\phantom{0}\phantom{0}n/a & \phantom{0}620 & yes & no \\
RoboMIND remainder & AgileX, humanoid & \phantom{0}\phantom{0}\phantom{0}n/a & \phantom{0}260 & yes & no \\
Kitchen and hand-object video & unconstrained & \phantom{0}\phantom{0}\phantom{0}n/a & \phantom{0}122 & no & partial \\
\midrule
\rowcolor{tabhl}
\textbf{Total} & & & \textbf{6\,451} & & \\
\bottomrule
\end{tabular}}
\end{table}

\subsection{Pretraining and Downstream Use}
\label{sec:procedure}

\Paragraph{Cached encoders.}
The frozen autoencoder and the frozen text encoder are evaluated offline. Latent
shards store the three latent frames of each window in bfloat16 after the
autoencoder's own standardisation, and text embeddings are deduplicated by
prompt so that a corpus with a few thousand distinct instructions stores a few
thousand embeddings. The cache key is the tuple of source, canvas, window,
normalisation, codec, and consequence layout, and any change to one element
rebuilds the cache. Training reads the cache, so the frozen modules stay outside
the loop, which removes the largest fixed cost from the step time and makes an arm of the
coupling axis cheap enough to run several times.

\Paragraph{Optimisation.}
Training uses fully sharded data parallelism \citep{zhao2023fsdp,paszke2019pytorch} with bfloat16
compute, single-precision master weights and optimiser state, and per-layer
activation recomputation. The optimiser is AdamW at a base learning rate of one
times ten to the minus four with momentum coefficients 0.9 and 0.95, weight decay
0.01, gradient clipping at unit norm, and a cosine schedule with linear warmup
down to one percent of the peak. The per-device batch holds sixteen windows.
Checkpoints are written atomically at a fixed step interval and the job chain
resumes from the newest one, which is what lets a long run survive the preemption
policy of the cluster. Seeds fix the data order, the noise draw, and the dropout
mask, and the evaluation seed is configured separately from the training seed.

\Paragraph{Downstream use.}
A published checkpoint is a self-contained directory holding the resolved
configuration, the merged weights, the component assets, the normalisation
statistics, and the data fingerprint. Deployment loads that directory alone and
serves a policy behind a websocket protocol that existing benchmark clients
speak without modification, which keeps the evaluation harness independent of
the training code. The policy walks the joint path of \sref{sec:paths}, returns
an action chunk in the platform's native units after inverting the layout and the
normalisation, and executes it under receding-horizon control. The other two
paths are driven from the checkpoint by the analysis code, so the representation
measurements read the weights directly.

\subsection{From the Formulation to the Measurements}
\label{sec:plan}

The formulation makes one claim testable: coupling action generation to world
generation changes measurable properties of the world representation, and those
changes rather than parameter count or supervision volume explain downstream
behaviour. Four questions decompose that claim, and each one is decided by an
instrument this section has already defined.

Whether a passive representation is close to action-free is decided by the
world-clamped ray against a randomly initialised backbone, and
\sref{sec:probes} reports it. Whether the gain comes from shaping the world
stream or from reading it is decided by the gradient degree of freedom of the
coupling axis, and \sref{sec:sweeps} reports it. Whether the generated future
follows an intervention on the action is decided by the action-clamped ray
against a shuffled chunk and a zero chunk, and \sref{sec:clamped} reports it.
Whether the gain concentrates outside the training distribution is decided by
the perturbation axes, the randomised split, and the held-out arm types, and
\sref{sec:closedloop} reports it. Every instrument carries a declared floor, and
a measurement that fails to clear its floor retracts the question it was meant
to decide.
